\begin{textsample}{POS Dim 1 – es – Score 39.00 – es/es\_ef\_anos\_iniciais\_ciencias\_da\_natureza\_2.txt}  \label{ex:f1_pos_199}
O Componente Curricular Ciências O ensino das Ciências da Natureza se concretiza no Ensino Fundamental através do componente curricular chamado simplesmente de “ Ciências ”.
Mas, tal simplicidade no nome não encerra toda a complexidade histórica, que vem a se desdobrar no Ensino Médio nos componentes curriculares de Biologia, Física e Química.
No Ensino Fundamental, o componente curricular “ Ciências ” carrega o privilégio e a responsabilidade de impulsionar a natural curiosidade dos estudantes pelo mundo que os cerca, estabelecendo as bases do pensamento científico e o prazer por continuar aprendendo ( FURMAN, 2009 ).
Diante dessa responsabilidade, a BNCC traz para a área de \textbf{Ciências} da Natureza um compromisso com o desenvolvimento do \textbf{letramento} científico : > Na escrita deste documento curricular, buscou -se a essência desse conceito, para além das discussões a respeito de qual o melhor termo a ser utilizado, seja \textbf{letramento}, alfabetização ou enculturação científica.
Buscou -se de fato demonstrar que a intenção do Ensino de \textbf{Ciências} é prover “ condições para que temas e situações envolvendo as \textbf{ciências} sejam analisados à luz dos conhecimentos científicos, sejam estes conceitos ou aspectos do próprio fazer científico ” ( SASSERON, 2015 ). > Dessa forma, o processo investigativo deve ser entendido como elemento \textbf{central} na formação dos estudantes, em um sentido mais amplo, e cujo desenvolvimento deve ser atrelado a situações didáticas planejadas ao longo de toda a Educação Básica, de modo a possibilitar que eles revisitem de forma \textbf{reflexiva} seus conhecimentos e sua \textbf{compreensão} acerca do mundo em que vivem. > ( BRASIL, 2017 ) > O \textbf{letramento} científico é entendido como instrumento que permite ao estudante o exercício pleno de sua cidadania, ao desenvolver sua \textbf{capacidade} de atuação no e sobre o mundo, permitindo -lhe compreendê -lo e interpretá -lo em seus aspectos natural, social e tecnológico, e transformá -lo com base nos aportes \textbf{teóricos} e processuais das Ciências. > ( BRASIL, 2017 ) > “... assim como a própria \textbf{ciência}, a Alfabetização Científica deve estar sempre em construção, englobando novos conhecimentos pela \textbf{análise} e em decorrência de novas situações ; de mesmo modo, são essas situações e esses novos conhecimentos que \textbf{impactam} os processos de construção de entendimento e de tomada de decisões e \textbf{posicionamentos} e que evidenciam as relações entre as \textbf{ciências}, a sociedade e as distintas áreas de conhecimento, ampliando os âmbitos e as perspectivas associadas à Alfabetização Científica. ” > ( SASSERON, 2015 ) Dessa forma, o ensino das Ciências da Natureza deve ser pensado como uma preparação para a vida, permitindo que o estudante não se \textbf{restrinja} à admiração passiva de uma “ \textbf{Ciência} dos \textbf{cientistas} ”, mas que transcenda o papel de espectador para o papel de criador de conhecimento ( MENEZES, 2009 ).
Ao mesmo tempo, deve -se evitar a concepção de que a \textbf{Ciência} seja algo \textbf{restrito} aos laboratórios, assumindo -se uma postura em que as práticas cotidianas sejam objeto de estudo e pesquisa ( BIZZO, 2002 ).
Partindo dessas concepções, a BNCC estabelece os seguintes pressupostos \textbf{teóricos} para o ensino de \textbf{Ciências} da Natureza ( BRASIL, 2017 ) : \textbf{Figura} 16 - Pressupostos \textbf{teóricos} para o ensino de \textbf{Ciências} na BNCC.
Cada um desses pressupostos \textbf{teóricos} possui práticas a serem desenvolvidas por estudantes e professores para que se concretizem na aquisição de novas habilidades e \textbf{competências} ( BRASIL, 2017 ) : \textbf{Figura} 17 - Práticas para os pressupostos \textbf{teóricos} “ Definição de \textbf{problemas} ” e “ Intervenção ”.
Ensino de Ciências - Definição de \textbf{problemas} - Levantamento, \textbf{análise} e representação - Comunicação - Intervenção Definição de \textbf{problemas} - Observar o mundo a sua volta e fazer perguntas. - Analisar \textbf{demandas}, delinear \textbf{problemas} e planejar investigações. - Propor hipóteses.
Intervenção - Implementar soluções e avaliar sua eficácia para resolver \textbf{problemas} cotidianos. - Desenvolver ações de intervenção para melhorar a qualidade de vida individual, coletiva e \textbf{socioambiental}.
Levantamento, \textbf{análise} e representação - Planejar e realizar atividades de campo ( experimentos, observações, leituras, visitas, ambientes \textbf{virtuais} etc. ). - Desenvolver e utilizar ferramentas, inclusive \textbf{digitais}, para coleta, \textbf{análise} e representação de dados em diferentes linguagens e formatos de mídias. - Avaliar informação ( validade, coerência e adequação ao \textbf{problema} formulado ). - Elaborar explicações e/ou modelos. - Associar explicações e/ou modelos à evolução histórica dos conhecimentos científicos envolvidos. - Selecionar e construir \textbf{argumentos} com base em evidências, modelos e/ou conhecimentos científicos. - \textbf{Aprimorar} seus saberes e incorporar, gradualmente, e de modo significativo, o conhecimento científico. - Desenvolver soluções para \textbf{problemas} cotidianos usando diferentes ferramentas, inclusive \textbf{digitais}.
Comunicação - Organizar e/ou extrapolar conclusões. - Relatar informações de forma oral, escrita ou multimodal. - Apresentar, de forma sistemática, dados e resultados de investigações. - Participar de discussões de caráter científico com colegas, professores, familiares e comunidade em geral. - Considerar contra-argumentos para rever processos investigativos e conclusões.
Com base nesses pressupostos, o Currículo de \textbf{Ciências} para o Ensino Fundamental busca uma maior aproximação com a perspectiva de uma Educação Científica, conforme descrita por Pedro Demo : > “... de um \textbf{lado} alfabetização, no sentido de propedêutica da vida em sociedade, no mundo do trabalho, no exercício da cidadania, no entendimento da realidade.
Noutro é formação permanente, porque nos acompanha pela vida afora, em especial em seu sentido autocrítico. ” > ( DEMO, 2017, p. 30 ) E, dessa forma, o processo de \textbf{ensino-aprendizagem} de Ciências centra -se no diálogo entre estudantes e professores, transformando os espaços e tempos de aprendizagem em espaços e tempos de interação entre os conhecimentos socioculturais de toda a comunidade escolar.
O saber científico torna -se, então, um instrumento ou ferramenta que, unido aos conhecimentos dos outros componentes curriculares e aos saberes populares, contribui para a \textbf{compreensão} e a solução dos \textbf{problemas} mais \textbf{imediatos} para a comunidade escolar, à medida que promove a tomada de consciência das possibilidades e dos limites das \textbf{competências} mediadoras de cada um ( ESPÍRITO SANTO, 2009 ).
Na perspectiva desse diálogo entre estudantes e professores, o domínio dos sistemas linguísticos populares e científicos torna -se essencial, pois são instrumentos socioculturais por meio dos quais os estudantes conhecem e compreendem as \textbf{complexas} interações dos conhecimentos que estão presentes nas suas práticas cotidianas e que, de alguma forma, explicam a condição humana.
Nesse sentido, tal domínio permite aos estudantes compreenderem não só a sua condição humana, como também as diferenças culturais \textbf{inerentes} a todo ser humano e à sociedade de modo geral ( MORIN, 1984 ).
Nessa concepção, compreender a diferença cultural significa, entre outras coisas, aceitar as diferentes formas de conhecer e explicar a condição humana, uma vez que a produção dos conhecimentos é sociohistórica, tornando todos os conhecimentos relativos e incertos.
Em consequência, o ensino de \textbf{Ciências} lida com essa incerteza dos saberes humanos, contribuindo para que cada estudante possa “ [... ] enfrentar as incertezas e, mais globalmente, o destino incerto de cada indivíduo e de toda a \textbf{humanidade} ” ( MORIN e LE MOIGNE, 2000, p. 55-56 ).
Vale relembrar que a BNCC define o desenvolvimento de uma \textbf{competência} como a \textbf{capacidade} de \textbf{mobilizar} conhecimentos ( conceitos e procedimentos ), habilidades ( práticas, cognitivas e \textbf{socioemocionais} ), atitudes e valores para resolver \textbf{demandas} \textbf{complexas} da vida cotidiana, do pleno exercício da cidadania e do mundo do trabalho.
E essa definição reflete uma concepção \textbf{teórica} que aborda a formação integral do indivíduo através do desenvolvimento de \textbf{competências}, ou seja, de conhecimentos, habilidades e atitudes.
Ela surge da necessidade de uma alternativa a modelos formativos, que priorizam o saber \textbf{teórico} sobre o prático, \textbf{sugerindo} uma \textbf{abordagem} pedagógica do conteúdo que possibilite aos estudantes um ensaio a respeito do saber conhecer, saber fazer e saber ser perante o objeto de estudo.
Nesta perspectiva, o desenvolvimento de \textbf{competências} no âmbito escolar “ deve identificar o que qualquer pessoa necessita para responder os \textbf{problemas} aos quais será exposta ao longo da vida ” ( ZABALA e ARNAU, 2010 ).
Tal \textbf{abordagem} também reflete as dimensões do conhecimento, conforme preconizadas pelos PCN ( BRASIL, 1997 ) e sua relação com os quatro pilares da Educação ( DELORS, 2012 ), quais sejam : | Dimensões do conhecimento | Pilares da Educação | |---|---| | Dimensão \textbf{conceitual} | Aprender a conhecer | | Dimensão procedimental | Aprender a fazer | | Dimensão atitudinal | Aprender a ser | | | Aprender a conviver | \textbf{Figura} 19 - Pilares da educação e as dimensões do conhecimento a serem desenvolvidas na \textbf{abordagem} por \textbf{competências}.
Considerando esses pressupostos \textbf{teóricos} e em articulação com as dez \textbf{competências} gerais e os saberes \textbf{socioemocionais}, a BNCC prevê o desenvolvimento de oito \textbf{Competências} Específicas para o ensino de Ciências da Natureza no Ensino Fundamental : Essa forma de organização reflete o que já ocorria nos \textbf{Parâmetros} Curriculares Nacionais e “ representam uma organização articulada de diferentes conceitos, procedimentos, atitudes e valores para cada um dos ciclos da escolaridade, compatível com os critérios de seleção acima apontados ” ( BRASIL, 1998 ).
Os conhecimentos e \textbf{tecnologias} desenvolvidos pelas diferentes \textbf{ciências} que compõem a área de \textbf{Ciências} da Natureza são um primeiro referencial para as \textbf{competências} específicas a serem desenvolvidas.
Dessa forma, as Unidades Temáticas são organizadas conforme \textbf{quadros} a seguir : Matéria e energia - Estudo de materiais e suas transformações - Fontes e tipos de energia utilizados na vida em geral ; - Ocorrência, utilização e processamento de recursos naturais e energéticos empregados na geração de diferentes tipos de energia ; - Produção e no uso responsável de materiais diversos ; - Perspectiva histórica da apropriação humana sobre os recursos naturais e energéticos ; - Identificação do uso de materiais em diferentes ambientes e épocas e sua relação com a sociedade e a \textbf{tecnologia}.
Vida e evolução - Características e necessidades dos seres vivos ( incluindo os seres humanos ) - A vida como fenômeno natural e social e elementos essenciais à sua manutenção ; - Processos evolutivos que geram a diversidade de formas de vida no planeta ; - Interações dos seres vivos com outros seres vivos e com os fatores não vivos do ambiente ; - Interações que os seres humanos estabelecem entre si e com os demais seres vivos e elementos não vivos do ambiente ; - Importância da preservação da biodiversidade e como ela se distribui nos principais ecossistemas brasileiros.
Terra e universo - Características da Terra, do Sol, da Lua e de outros corpos celestes ( dimensão, composição, localização, movimentos e forças atuantes ) ; - Experiências de observação do céu e dos principais fenômenos celestes ; - Observação do planeta Terra, particularmente das zonas habitadas pelo ser humano e demais seres vivos ; - Processo histórico de construção dos conhecimentos sobre a Terra e o céu em diferentes culturas ao longo da história da \textbf{humanidade} ; - Fenômenos naturais como vulcões, tsunamis e terremotos ; - Fenômenos naturais relacionados aos padrões de circulação atmosférica e oceânica e aos movimentos da Terra ; - Conhecimentos relativos à evolução da vida e do planeta, ao clima e à previsão do tempo, entre outros fenômenos. \textbf{Figura} 21 - Principais objetos de estudo das Unidades Temáticas de \textbf{Ciências} da Natureza.
Como se pode observar, os estudantes deverão desenvolver, ao longo do Ensino Fundamental, sua própria forma de ver e entender o mundo à sua volta, a partir de vivências em práticas investigativas, da apropriação de linguagens próprias das \textbf{ciências} e do estabelecimento de relações entre a \textbf{Ciência}, a \textbf{Tecnologia}, a Sociedade e o Ambiente em que estão inseridos ( BRASIL, 2016 ).
Nesse ponto, vale ressaltar o risco de se resvalar para dois limites que devem ser evitados, o de se ensinar uma “ \textbf{ciência} infantiloide e apequenada ” para os pequenos ou de se forçar uma proposta de Ciências que as crianças não teriam maturidade para levar adiante ( LINN e EYLON, 2011 ).
Diante disso, caberá às escolas e aos docentes uma revisão constante de suas práticas, de forma a buscar a “ dose \textbf{certa} de pretensão científica cabível ” ao ensino e à aprendizagem de \textbf{Ciências} da Natureza ( DEMO, 2017 ).

% matched lemmas: abordagem, análise, aprimorar, argumento, capacidade, central, certo, cientista, ciência, competência, complexo, compreensão, conceitual, demanda, digital, ensino-aprendizagem, figura, figurar, humanidade, imediato, impactar, inerente, lado, letramento, mobilizar, parâmetro, posicionamento, problema, quadro, reflexivo, restringir, socioambiental, socioemocional, sugerir, tecnologia, teórico, virtual
\end{textsample}
